\documentclass[12pt]{article}

\usepackage[margin=1in]{geometry}
\usepackage[utf8]{inputenc}
\usepackage[T1]{fontenc}
\usepackage{lmodern}
\usepackage{microtype}
\usepackage{setspace}
\usepackage{parskip}
\usepackage{url}
\usepackage{booktabs}
\usepackage{threeparttable}
\usepackage[font={footnotesize}]{caption}
\usepackage{float}
\usepackage{graphicx}
\usepackage{lscape}
\usepackage[super,comma,sort&compress]{natbib}
\bibliographystyle{unsrtnat}
\usepackage[colorlinks=true, urlcolor=blue, linkcolor=black, citecolor=black]{hyperref}

% Auto-generated by code/generate_faq_numbers.R. Do not edit by hand.
\newcommand{\faqNovelty}{1}
\newcommand{\faqNegotiatorPay}{2}
\newcommand{\faqFines}{3}
\newcommand{\faqProvision}{4}
\newcommand{\faqInsurance}{5}
\newcommand{\faqPharma}{6}
\newcommand{\faqPbmComparison}{7}
\newcommand{\faqComprehensiveScope}{8}
\newcommand{\faqDispensingFee}{9}
\newcommand{\faqPriceFixing}{10}
\newcommand{\faqCentralPlanning}{11}
\newcommand{\faqAllPayer}{12}
\newcommand{\faqInvertedSchedule}{13}
\newcommand{\faqBundles}{14}
\newcommand{\faqOptOut}{15}
\newcommand{\faqRural}{16}
\newcommand{\faqWithdrawal}{17}
\newcommand{\faqInnovation}{18}
\newcommand{\faqFiscal}{19}


% Allow figures and tables to float more freely: let them take up most of a page
% if needed, and do not leave large gaps chasing a perfect placement.
\renewcommand{\topfraction}{0.9}
\renewcommand{\bottomfraction}{0.8}
\renewcommand{\textfraction}{0.07}
\renewcommand{\floatpagefraction}{0.7}
\setcounter{topnumber}{3}
\setcounter{bottomnumber}{3}
\setcounter{totalnumber}{5}

\onehalfspacing

\title{One Drug, Many Prices:\\A Proposal to Reform US Healthcare Pricing}
\author{Maysam Rabbani\thanks{Department of Economics, Montclair State University. Email: \texttt{rabbanim@montclair.edu}.}}
\date{}

\begin{document}

\maketitle
\thispagestyle{empty}

\subsection*{Introduction}

The US healthcare pricing system is broken: the same product is sold to different payers for wildly different prices, pharmacy benefit managers collect undisclosed sums and hide the amounts as a tightly guarded secret\citep{ftc2024pbm}, and the largest public insurer is effectively barred from using its leverage to reduce prices. Americans usually do not know the price at the time of receiving care. This opacity is a structural driver of the price dispersion. No other developed nation operates a pricing system this opaque.\citep{anderson2003prices}

I propose three reforms: consolidated government negotiation, Bloc-wide pricing as a condition of market access, and a central price lookup system. Each has been tried or proposed in isolation with promising results. I show how the three could be combined into a comprehensive plan (see \href{https://rabbanimaysam.github.io/health-reform-proposal/faq.html}{FAQ \faqNovelty}). This note is an exercise in policy design. It presents an idea to invite consideration and critique. Legislative strategy and political roadmap would be next steps that are beyond the purview of this note.

Figure~\ref{fig:pricegaps} shows the payment gap across ten common drugs, hospital services, and outpatient procedures. For example, a colonoscopy costs \$880 at the VA, \$1,760 under Medicare, \$5,816 under a commercial insurer, and \$19,206 for uninsured patients. This is not an anomaly. It is common for US healthcare prices to swing several-fold across payers, and the uninsured, a vulnerable group, pay the most.

\begin{figure}[htbp]
\centering
\includegraphics[width=0.9\textwidth]{fig_pricegaps.pdf}
\caption{Price gaps across payers for the same service or product.}
\label{fig:pricegaps}
\vspace{2pt}
\begin{minipage}{\textwidth}
\fontsize{8}{9}\selectfont
\textit{Notes:} The "private insurer" bar is the total negotiated rate, including patient cost-sharing; the "uninsured" bar is the chargemaster or list price. Each panel uses its own vertical scale. Missing bars indicate that a reliable figure was not available in the cited sources. Full sources and derivations are on the \href{https://rabbanimaysam.github.io/health-reform-proposal/price-gaps.html}{companion site}.
\end{minipage}
\end{figure}

Two forces determine an insurer's ability to reduce prices: its size, and the availability of alternatives. A large insurer has more leverage because losing its volume hurts providers more, but if providers can walk to another insurer offering higher rates, that leverage evaporates. Leverage is maximized when the insurer is large and irreplaceable. The three reforms below are designed around this logic. Consolidated negotiation makes the government a single, large insurer. Bloc-wide pricing makes it difficult for providers to replace the Bloc. A central price lookup system ensures the negotiated prices are honored, not undercut by secret side deals.

\subsection*{Reform 1: Consolidated government negotiation}

Reform 1 pools all federal healthcare purchasing into a single negotiating Bloc that bargains on behalf of every federal program.

Today, public purchasing power is fragmented twice over. At the federal level, Medicare, Medicaid, and at least nine other programs each maintain their own formulary, fee schedule, and leverage. At the state level, each of the 50 states administers its own Medicaid program and its own employee health plans, adding another layer of fragmentation. For prescription drugs alone, federal and state programs combined account for roughly half of US drug spending, yet they compete against one another instead of bargaining as one.\citep{cms2024nhe} Hospital and physician payments are just as fragmented. Consolidating these programs into a single negotiating Bloc would use leverage the government has but wastes. Other developed nations leverage their full size in what resembles a bloc, and pay far less.\citep{gusmano2020price}

From Part D's launch in 2006 until 2022, Medicare was statutorily prohibited from negotiating drug prices, a restriction unparalleled in the developed world. The Inflation Reduction Act of 2022 began to reverse it, authorizing Medicare to negotiate prices for an initial ten drugs, with more added annually. Confining negotiation to ten drugs is an arbitrary impediment that forces the government to underutilize its leverage and wastes taxpayer money. The VA, unencumbered by that prohibition, pays half of Medicare Part D despite being much smaller.\citep{cbo2022commercial, cbo2021drugprices, gao2021vadrugs}

Figure~\ref{fig:intlprices} compares prices in the US and six other developed countries for nine of those ten drugs. The US list price towers over what every peer nation pays. Even after rebates and discounts, private insurers in the US pay multiples of what patients in Canada, France, or the UK pay for the same drug. Medicare's newly negotiated prices are 22\% lower\citep{wouters2025jama} and narrow the gap but remain well above other nations.

\begin{figure}[htbp]
\centering
\includegraphics[width=\textwidth]{fig_intlprices.pdf}
\caption{Drug prices by country.}
\label{fig:intlprices}
\vspace{2pt}
\begin{minipage}{\textwidth}
\fontsize{8}{9}\selectfont
\textit{Notes:} The solid dark bar is the US list price (paid by the uninsured). The striped bar is the estimated US net price across all US payers combined, after rebates and discounts.\citep{wouters2025jama} The crosshatched bar is the Medicare negotiated price under the Inflation Reduction Act (effective 2026). Insulin aspart is excluded due to reporting incompatibility across countries. Each panel uses its own vertical scale. Full figures and sources are on the \href{https://rabbanimaysam.github.io/health-reform-proposal/international-prices.html}{companion site}.
\end{minipage}
\end{figure}

Reform 1 outlaws rebates, chargebacks, spread pricing, and side payments of any kind. Negotiation is conducted by salaried Bloc negotiators with expertise in pharmacology, clinical medicine, and health economics, who also make formulary and tiering decisions. Bloc negotiators are prohibited from accepting any payment beyond their government compensation package, and for twelve months after leaving the Bloc are barred from employment in any sector of the regulated healthcare industry. Compensation follows the federal Senior Executive Service model: a fixed salary at the SES level and a performance bonus capped at 10\% of base, awarded on periodic peer-reviewed evaluations rather than per-contract outcomes (see \href{https://rabbanimaysam.github.io/health-reform-proposal/faq.html}{FAQ \faqNegotiatorPay}).

\textcolor{blue}{Whereas} PBMs aggressively negotiate price reductions, a centralized Bloc-wide negotiation would be advantageous on four grounds: it eliminates the leakage that results from the secretive rebate system; instead of limiting the discount to a subset of patients and leaving others exposed to list prices, it would extend the discounted price to everyone; it eliminates the inherent conflict of interest that arises from the fact that PBMs are largely owned by health insurance companies and retail pharmacy chains;\citep{ftc2024pbm} and the difficulty of substituting away from the Bloc equips it with substantially more negotiation power than PBMs, which could lead to significantly lower prices. See \href{https://rabbanimaysam.github.io/health-reform-proposal/faq.html}{FAQ \faqPbmComparison} for the detailed \textcolor{green}{comparison}.

Smaller-scale precedents for a bloc negotiation exist. The Federal Supply Schedule publishes VA-negotiated contract prices that the Department of Defense (including Tricare), the Indian Health Service, the Federal Bureau of Prisons, and the Public Health Service all purchase off,\citep{vafss2026} demonstrating that one body can negotiate for many federal agencies. At the state level, pooled purchasing coalitions such as the Sovereign States Drug Consortium extract supplemental Medicaid rebates beyond federal levels by combining state purchasing power.\citep{oig2014supplementalrebates, horvath2019purchasing}

\textcolor{blue}{Reform} 1 grants the Bloc authority over all drugs and services without exclusions or carve-outs, covering even products and services that are reasonably priced today, such as generics. Universal scope is not universal action: where competition produces low prices on its own, the Bloc need not negotiate. But authority that stops short of universal leaves doors open for exploitation, and every carve-out becomes a roadmap for the next round of pricing abuse (see \href{https://rabbanimaysam.github.io/health-reform-proposal/faq.html}{FAQ \faqComprehensiveScope}).

\subsection*{Reform 2: Bloc-wide pricing as a condition of market access}

Any provider that contracts with the Bloc retains the freedom to choose which private insurers to contract with, but the contracted price must be the same as the Bloc-negotiated rate. Uninsured patients cannot be refused or charged a different price. The negotiated price is locked in: no insurer can negotiate further discounts on the side. Any provider that violates these terms loses its contract with the Bloc, and with it, access to the majority of the market.

The cost problem largely lies with hospitals and pharmaceuticals, not physicians. Hospital prices for the privately insured grew 42\% between 2007 and 2014, while physician prices for the same services grew only 18\%.\citep{cooper2019hospitalprice, cooper2019price, gowrisankaran2015mergers, dafny2009mergers} Yet physicians must be included in the Bloc-wide pricing. Hospitals routinely acquire physician practices and, once they do, prices rise.\citep{cooper2025physicianacquisitions} Leaving physicians outside the Bloc would create a back door for hospitals to shift billable activity to nominally independent physician groups, circumventing the Bloc's pricing rules.

Reform 2 does not make insurance plans identical; it standardizes the price, not the coverage. Once the price of a service is fixed, insurers would compete on what fraction of that price they cover and what fraction the patient pays out of pocket. Generous plans leave the patient with a small copay, cheaper plans with more, but both pay the provider the same.

\textcolor{blue}{Price} discrimination is theoretically ambiguous and can produce equity gains when the more elastic buyer is the lower-income one.\citep{varian1985pd} The case against it in US healthcare is empirical, not theoretical.\citep{reinhardt2011allpayer, anderson2015allpayer} The standard sorting is inverted. Price tracks leverage, and leverage is thin in many transactions. The uninsured face list prices.\citep{dusetzina2016uninsured, bai2015markup} Insured patients pay the full negotiated rate at the point of care until the deductible is met, functioning as uninsured for a meaningful fraction of the year. Specialty-drug patients on coinsurance pay a percentage of list, not of the post-rebate net. Out-of-network and emergency patients face near-list rates. Cash payers using third-party discount tools sometimes pay less for the same drug at the same pharmacy than their own insurance plan would charge. Small employers and small commercial insurers pay higher rates than large ones for the same care.

The inversion is the operating principle of the schedule, not a failure at one tail of the distribution. Hospital market power intensifies the harm wherever leverage is thin,\citep{cooper2019price} and price opacity shields the schedule from scrutiny. Solving for the uninsured alone would protect the most visible group while leaving the structure intact for the latent ones. Reform 2 replaces this schedule with uniform \textcolor{green}{pricing} (see \href{https://rabbanimaysam.github.io/health-reform-proposal/faq.html}{FAQ \faqInvertedSchedule}).

Existing attempts at constraining price discrimination have proven mostly successful. Maryland has run an all-payer hospital system since 1977.\citep{murray2009maryland} The 340B Program requires Medicaid manufacturers to sell to safety-net providers at a statutory discount (at least 22.5\% off average sales price),\citep{medpac2015340b} though the program has left room for exploitation.\citep{conti2014340b} The Medicaid Best Price rule requires manufacturers to offer Medicaid the lowest price offered to any non-federal purchaser (see \href{https://rabbanimaysam.github.io/health-reform-proposal/faq.html}{FAQ \faqAllPayer} for why this is not all-payer pricing).

\subsection*{Reform 3: A central price lookup system}

Reforms 1 and 2 produce one negotiated price per provider per service. Reform 3 makes that price visible. The Bloc operates a public price lookup system, the canonical record of every contracted price in the country. Each entry pairs a service code with a provider and a price. It updates whenever a contract is signed or revised.

Every Bloc contract pre-sets the price of every covered service, whether per service or per bundled episode. Bloc officials directly transfer prices from contract to database, with no provider input. This sidesteps the misreporting, gaming, and non-compliance that become systemic whenever transparency depends on thousands of providers posting their own data.

The lookup system is open, free, and built for comparison. A patient may compare hip replacement cost in every hospital within a chosen radius, or see EpiPen prices at every nearby pharmacy in real time. A public API lets third parties integrate prices into comparison tools, electronic health records, and patient portals.

Physicians can use the same tool. A doctor writing a prescription for an expensive drug can identify the pharmacy with the lowest local price and send the prescription there. For costly procedures, patients may be referred across state lines if the savings justify the trip. This intensifies competition: physicians, better informed and more influential than patients in directing demand, become active routers of patients to lower-priced sellers.

Price transparency is a known driver of price reductions,\citep{robinson2013calpers, whaley2019nh} but only if the posted prices are the actual prices and if transparency is enforced (see \href{https://rabbanimaysam.github.io/health-reform-proposal/faq.html}{FAQ \faqNovelty}). Existing federal transparency rules have failed on both counts:\citep{jiang2022transparency} they require posting in formats too inconsistent to compare, and punish non-compliance with symbolic penalties. Under Reform 3, the Bloc posts actual prices, and overcharged patients can file a complaint. Deterrence requires large fines, for example a multiplier on the order of 100 times the overcharge (see \href{https://rabbanimaysam.github.io/health-reform-proposal/faq.html}{FAQ \faqFines}), and contract termination for repeat offenders.

\subsection*{Final word}

The reforms would benefit patients most. They would pay less, comparison-shop, and know prices before receiving care. As underlying prices fall, insured patients would face lower premiums, and uninsured patients would no longer be overcharged. The US has over-spent on healthcare for decades. The tools to change that are not exotic. They are familiar, tested at smaller scales, and in partial federal use. What is missing is the decision to put them together.

The \href{https://rabbanimaysam.github.io/health-reform-proposal/faq.html}{online FAQ} answers a set of common questions about the proposal. Comments, criticisms, and suggestions are welcome through \href{https://rabbanimaysam.github.io/health-reform-proposal/}{this link}.

{\scriptsize
\setlength{\bibsep}{0pt}
\bibliography{references}
}

\end{document}
